\section*{Exercises on \S\arabic{exercisegroup}}
\phantomsection
\addcontentsline{toc}{section}{Exercises on \protect\S\arabic{exercisegroup}}

\begin{enumerate}
\def\labelenumi{\arabic{enumi})}
\tightlist
\item
  \begin{enumerate}
  \def\labelenumii{\alph{enumii})}
  \tightlist
  \item
    With the notation of IV, p.~165, suppose that \(\mathbf{f}\) is uniformly continuous on \(\mathrm{J} \times \mathrm{S}\) . Prove prop. 3 of IV, p.~166 without using the axiom of choice. (Let \(\delta\) be such that the relations \(|t_2 - t_1| \leqslant \delta\) , \(\| \mathbf{x}_2 - \mathbf{x}_1 \| \leqslant \delta\) imply \(\| \mathbf{f}(t_2, \mathbf{x}_2) - \mathbf{f}(t_1, \mathbf{x}_1) \| \leqslant \varepsilon\) ; consider a subdivision of J into intervals of length \(\leqslant \inf(\delta, \delta / M)\) and define the approximate solution by induction.)
  \end{enumerate}
\end{enumerate}

\begin{enumerate}
\def\labelenumi{\alph{enumi})}
\setcounter{enumi}{1}
\tightlist
\item
  When E is finite dimensional and f is Lipschitz on \(I \times H\) prove prop. 3 of IV, p.~166 without using the axiom of choice. (Remark that for every \(\delta > 0\) there exists a finite number of points \(x_{i}\) of S such that every point of S is at a distance \(\leqslant \delta\) from one of the \(x_{i}\) ; then proceed as in a), considering the regulated functions \(t \mapsto \mathbf{f}(t, \mathbf{x}_{i})\) , of which there are only a finite number.)
\end{enumerate}

\begin{enumerate}
\def\labelenumi{\arabic{enumi})}
\setcounter{enumi}{1}
\item
  Given two numbers \(b > 0\) , \(M > 0\) and an arbitrary number \(\varepsilon > 0\) , give an example of a scalar differential equation \(x' = g(x)\) such that \(|g(x)| \leqslant M\) for \(|x| \leqslant b\) and which admits an integral \(x = u(t)\) continuous on the interval \(] - b/M - \varepsilon, b/M + \varepsilon[\) , but which does not have a finite limit at the point \(x = \frac{b}{M} + \varepsilon\) (define \(g\) so that the integral considered has a continuous derivative on \(] - b/M - \varepsilon, b/M + \varepsilon[\) , this derivative being equal to the constant \(M\) on \(] - b/M, b/M[\) ).
\item
  Let \(S \subset H\) be an open ball with centre \(\mathbf{x}_0\) and radius \(r\) , and \(\mathbf{f}\) a Lipschitz function for the constant \(k > 0\) on \(I \times S\) ; suppose further that \(t \mapsto \mathbf{f}(t, \mathbf{x}_0)\) is bounded on \(I\) , and denote by \(M_0\) the supremum of \(\| \mathbf{f}(t, \mathbf{x}_0) \|\) over \(t \in I\) . Show that for all \(t_0 \in I\) there is an integral \(\mathbf{u}\) of \(\mathbf{x}' = \mathbf{f}(t, \mathbf{x})\) , with values in \(S\) , equal to \(\mathbf{x}_0\) at the point \(t_0\) and defined on the intersection of \(I\) and the interval \(]t_0 - \lambda, t_0 + \lambda[\) , where
\end{enumerate}

\[
\lambda = \frac {1}{k} \log \left(1 + \frac {k z}{M _ {0}}\right)
\]

(note that one has \(\| \mathbf{u}(t) - \mathbf{x}_0\| \leqslant \mathbf{M}(t - t_0) + k\int_{t_0}^t\| \mathbf{u}(s) - \mathbf{x}_0\| ds\) for \(t > t_0\) ).

* 4) a) Let \(I = ]t_0 - a, t_0 + a[\) be an open interval in \(\mathbf{R}\) , let \(S\) be an open ball with centre \(\mathbf{x}_0\) and radius \(r\) in \(E\) , and \(\mathbf{f}\) a locally Lipschitz function on \(I \times S\) . Let \((s, z) \mapsto h(s, z)\) be a function \(\geqslant 0\) defined and continuous for \(0 \leqslant s \leqslant a\) and \(0 \leqslant z \leqslant r\) , such that for every \(s \in [0, a]\) the map \(z \mapsto h(s, z)\) is increasing; suppose that \(\| \mathbf{f}(t, \mathbf{x}) \| \leqslant h(|t - t_0|, \| \mathbf{x} - \mathbf{x}_0 \|)\) on \(I \times S\) . Let \(\varphi\) be a primitive of a regulated function on an interval \([0, \alpha]\) (with \(\alpha < a\) ), with values in \([0, r[\) , such that \(\varphi(0) = 0\) and \(\varphi'(s) > h(s, \varphi(s))\) on \([0, \alpha]\) except on a countable set. Show that the integral \(\mathbf{u}\) of \(\mathbf{x}' = \mathbf{f}(t, \mathbf{x})\) , equal to \(\mathbf{x}_0\) at the point \(t_0\) , is defined on \(]t_0 - \alpha, t_0 + \alpha[\) , and that \(\| \mathbf{u}(t) - \mathbf{x}_0 \| \leqslant \varphi(|t - t_0|)\) on this interval.

\begin{enumerate}
\def\labelenumi{\alph{enumi})}
\setcounter{enumi}{1}
\tightlist
\item
  Deduce from a) that if f is defined on I × E, and there is a function h defined, continuous, increasing and \textgreater{} 0 on {[}0, +∞{[}, such that \(\int_{0}^{+\infty} dz / h(z) = +\infty\) , and that \(\|\mathbf{f}(t, \mathbf{x})\| \leqslant h(\|x\|)\) , then every integral of \(\mathbf{x}' = \mathbf{f}(t, \mathbf{x})\) is defined on all of I.
\end{enumerate}

* 5) Let \(E\) be the complete normed space of sequences \(\mathbf{x} = (x_{n})\) of real numbers such that \(\lim_{n\to \infty}x_n = 0\) , normed by \(\| \mathbf{x}\| = \sup_n|x_n|\) . For every integer \(n\geqslant 0\) let \(\mathbf{e}_n\) be the sequence every term of which is zero, apart from the term with index \(n\) , which is equal to 1; the space \(E\) is the direct sum of the subspace \(V_{n}\) of dimension 2 generated by \(\mathbf{e}_n\) and \(\mathbf{e}_{n + 1}\) , and the closed subspace \(W_{n}\) generated by the \(\mathbf{e}_k\) with indices different from \(n\) and \(n + 1\) . Let \(\mathbf{f}_n\) be a continuous Lipschitz function on \(E\) , with values in \(E\) , constant on every equivalence class modulo \(W_{n}\) , equal to \(\mathbf{e}_{n + 1} - \mathbf{e}_n\) on the line joining \(\mathbf{e}_n\) to \(\mathbf{e}_{n + 1}\) , and equal to 0 on the intersection of \(V_{n}\) and the ball \(\| \mathbf{x}\| \leqslant \frac{1}{4}\) . On the other hand, for every integer \(n > 0\) let \(\varphi_{n}\) be a real function defined and continuous on the interval \(\left[\frac{1}{n + 1},\frac{1}{n}\right]\) , equal to 0 at the endpoints of this interval, and such that \(\int_{1 / n + 1}^{1 / n}\varphi_n(t)dt = 1\) . Let \(I\) be the interval {[}0, 1{]} in \(\mathbf{R}\) ; consider on \(I\times E\) the function \(\mathbf{f}\) with values in \(E\) , defined as follows: \(\mathbf{f}(0,\mathbf{x}) = 0\) for every \(\mathbf{x}\in E\) ; for \(\frac{1}{n + 1}\leqslant t\leqslant \frac{1}{n}\) , let \(\mathbf{f}(t,\mathbf{x}) = \varphi_n(t)\mathbf{f}_n(\mathbf{x})\) . Show that \(\mathbf{f}\) is continuous and locally Lipschitz on \(I\times E\) , but that there exists an integral \(\mathbf{u}\) of the equation \(\mathbf{x}' = \mathbf{f}(t,\mathbf{x})\) , defined and bounded on {[}0, 1{]}, equal to \(\mathbf{e}_n\) for \(t = 1 / n\) , and consequently having no limit as \(t\) tends to 0.

* 6) a) Let \(\mathrm{I} = [0, +\infty[\) ; if \(\mathbf{f}\) is Lipschitz on \(\mathrm{I} \times \mathrm{E}\) for a regulated function \(k(t) > 0\) such that the integral \(\int_0^{+\infty} k(t) dt\) converges, show that every integral of the equation \(\mathbf{x}' = \mathbf{f}(t, \mathbf{x})\) is defined on all of \(\mathrm{I}\) .

\begin{enumerate}
\def\labelenumi{\alph{enumi})}
\setcounter{enumi}{1}
\tightlist
\item
  If further the integral \(\int_{0}^{+\infty}\|\mathbf{f}(t,\mathbf{x}_{0})\|dt\) converges (for a certain point \(x_{0}\in E\) ), show that every integral of \(\mathbf{x}^{\prime}=\mathbf{f}(t,\mathbf{x})\) tends to a finite limit as t tends to \(+\infty\) (first show that every integral is bounded as t tends to \(+\infty\) ).
\end{enumerate}

\begin{enumerate}
\def\labelenumi{\arabic{enumi})}
\setcounter{enumi}{6}
\tightlist
\item
  Consider the system of scalar differential equations
\end{enumerate}

\[
\frac {d x _ {i}}{d t} = \sum_ {j, k = 1} ^ {n} c _ {i j k} x _ {j} x _ {k} \quad (1 \leqslant i \leqslant n)
\]

where the \(c_{ijk}\) are constants such that \(c_{kji} = -c_{ijk}\) . Show that the integrals of this system are defined for all values of t (note that \(\sum_{i=1}^{n} x_{i}^{2}\) is constant for every integral \(\mathbf{x} = (x_{i})\) ).

\begin{enumerate}
\def\labelenumi{\arabic{enumi})}
\setcounter{enumi}{7}
\tightlist
\item
  Let \(\mathbf{f}\) be a function defined on \(\mathrm{I} \times \mathrm{H}\) satisfying the conditions of lemma 1 of IV, p.~165, and such that for a constant \(k\) with \(0 < k < 1\) , and a point \(t_0 \in \mathrm{I}\) , one has, for all \(t \in \mathrm{I}\) and \(\mathbf{x}_1\) , \(\mathbf{x}_2 \in \mathrm{H}\) ,
\end{enumerate}

\[
\left\| \mathbf {f} (t, \mathbf {x} _ {1}) - \mathbf {f} (t, \mathbf {x} _ {2}) \right\| \leqslant \frac {k}{\left| t - t _ {0} \right|} \left\| \mathbf {x} _ {1} - \mathbf {x} _ {2} \right\|.
\]

Show that if \(\mathbf{u}\) and \(\mathbf{v}\) are two approximate solutions of the equation \(\mathbf{x}' = \mathbf{f}(t, \mathbf{x})\) , to within \(\varepsilon_1\) and \(\varepsilon_2\) respectively, defined on I, with values in H, and having the same value at the point \(t_0\) , then, for every \(t \in I\) ,

\[
\left\| \mathbf {u} (t) - \mathbf {v} (t) \right\| \leqslant \frac {\varepsilon_ {1} - \varepsilon_ {2}}{1 - k} | t - t _ {0} |.
\]

Deduce from this that th. 1 (IV, p.~171) and 2 (IV, p.~172) remain valid for the equation \(\mathbf{x}' = \mathbf{f}(t, \mathbf{x})\) under the indicated conditions.

\begin{enumerate}
\def\labelenumi{* \arabic{enumi})}
\setcounter{enumi}{8}
\tightlist
\item
  Let I be an interval in \(\mathbf{R}\) , and \(t_0\) a point of I, let S be a ball with radius \(r\) and centre \(\mathbf{x}_0\) in E, let G be the normed space of bounded maps from \(\mathrm{I} \times \mathrm{S}\) into E, the norm of such a map \(\mathbf{f}\) being the supremum \(\| \mathbf{f} \|\) of \(\| \mathbf{f}(t, \mathbf{x}) \|\) over \(\mathrm{I} \times \mathrm{S}\) . For every \(M > 0\) let \(\mathrm{G}_{\mathrm{M}}\) be the ball \(\| \mathbf{f} \| \leqslant \mathrm{M}\) in G. Let L be the subset of G formed by the Lipschitz maps of \(\mathrm{I} \times \mathrm{S}\) into E; for every function \(\mathbf{f} \in \mathrm{L} \cap \mathrm{G}_{\mathrm{M}}\) let \(\mathbf{u} = \mathrm{U}(\mathbf{f})\) be the integral of \(\mathbf{x}' = \mathbf{f}(t, \mathbf{x})\) such that \(\mathbf{u}(t_0) = \mathbf{x}_0\) , with values in S, and defined on the intersection \(\mathrm{J}_{\mathrm{M}}\) of I and the interval
\end{enumerate}

\[
\left] t _ {0} - \frac {r}{\mathbf {M}}, t _ {0} + \frac {r}{\mathbf {M}} \right[ \text {(th. 1).}
\]

\begin{enumerate}
\def\labelenumi{\alph{enumi})}
\tightlist
\item
  Let \((\mathbf{f}_{n})\) be a sequence of functions belonging to \(L \cap G_{M}\) ; if \(f_{n}\) converges uniformly on \(I \times S\) to a function f, then every cluster point (for the topology of compact convergence) of the sequence of \(\mathbf{u}_{n} = \mathrm{U}(\mathbf{f}_{n})\) in the space F of bounded maps from \(J_{M}\) into E, is an integral of \(\mathbf{x}' = \mathbf{f}(t, \mathbf{x})\) taking the value \(x_{0}\) at the point \(t_{0}\) . Conversely, every integral v of \(\mathbf{x}' = \mathbf{f}(t, \mathbf{x})\) defined on \(J_{M}\) and such that \(\mathbf{v}(t_{0}) = \mathbf{x}_{0}\) , is also an integral of an equation \(\mathbf{x}' = \mathbf{g}(t, \mathbf{x})\) , where g is Lipschitz and arbitrarily close to f in G (consider the equation
\end{enumerate}

\[
\mathbf {x} ^ {\prime} = \mathbf {f} _ {n} (t, \mathbf {x}) + \mathbf {v} ^ {\prime} (t) - \mathbf {f} _ {n} (t, \mathbf {v} (t))).
\]

\begin{enumerate}
\def\labelenumi{\alph{enumi})}
\setcounter{enumi}{1}
\tightlist
\item
  Suppose further that E is of finite dimension. Show that if \(f \in G_{M}\) satisfies the hypotheses of lemma 1, then for every \(t \in J_{M}\) the set \(H(t)\) of values at the point t of the integrals of \(\mathbf{x}' = \mathbf{f}(t, \mathbf{x})\) which take the value \(x_{0}\) at the point \(t_{0}\) is a compact and connected set (to see that \(H(t)\) is closed, use Ascoli's theorem; to see that it is connected, use a): if \(x_{1}, x_{2}\) are two points of \(H(t)\) and \(\varepsilon > 0\) is arbitrary, show that there exists a connected set \(\Phi\) of functions g belonging to \(L \cap G_{M}\) such that \(\|f - g\| \leqslant \varepsilon\) for every function \(g \in \Phi\) , and that the set of values of the functions \(U(\mathbf{g})\) at the point t is connected and contains \(x_{1}\) and \(x_{2}\) . Conclude by passing to the limit along an ultrafilter finer than the neighbourhood filter of 0 in \(R_{+}\) ).
\end{enumerate}

\begin{enumerate}
\def\labelenumi{\arabic{enumi})}
\setcounter{enumi}{9}
\tightlist
\item
  Let \(f\) be a continuous bounded real function defined on the box \(\mathrm{P}: |t - t_0| < a\) , \(|x - x_0| < b\) of \(\mathbf{R}^2\) . Let \(\mathbf{M}\) be the supremum of \(|f(t,x)|\) over \(\mathrm{P}\) , and let \(\mathrm{I} = ]t_0 - \alpha, t_0 + \alpha[\) , where \(\alpha = \inf(a, b/\mathrm{M})\) . Show that the upper and lower envelopes of the set \(\Phi\) of integrals of \(x' = f(t,x)\) defined on \(\mathrm{I}\) and taking the value \(x_0\) at the point \(t_0\) , are again integrals of \(x' = f(t,x)\) , which one calls the maximal and minimal integral of this equation corresponding to the point \((t_0, x_0)\) (remark that the set \(\Phi\) is equicontinuous and closed for the topology of uniform convergence on \(\mathrm{I}\) ).
\end{enumerate}

For every \(\tau \in I\) let \(\xi\) be the value of the minimal integral (corresponding to \((t_0, x_0)\) ) at the point \(\tau\) . Show that the minimal integral corresponding to the point \((\tau, \xi)\) is identical to the minimal integral corresponding to the point \((t_0, x_0)\) on an interval of the form \([\tau, \tau + h[\) if \(\tau > t_0\) , and of the form \(]\tau - h, \tau]\) if \(\tau < t_0\) .

Deduce from this that there exists a largest open interval \(]t_{1}, t_{2}[\) containing \(t_{0}\) and contained in \(]t_{0}-a, t_{0}+a[\) , such that the minimal integral u corresponding to \((t_{0}, x_{0})\) can be extended by continuity to \(]t_{1}, t_{2}[\) , so that at every point t of \(]t_{1}, t_{2}[\) , the value \(u(t)\) belongs to \(]x_{0}-b, x_{0}+b[\) , and that u is identical with the minimal integral corresponding to the point \((t, u(t))\) on an interval of the form \([t, t+h[\) if \(t > t_{0}\) and of the form \(]t-h, t]\) if \(t < t_{0}\) ; show further that either \(t_{1} = t_{0} - a\) (resp. \(t_{2} = t_{0} + a\) ) or that \(\lim_{t \to t_{1}} u(t) = x_{0} \pm b\) (resp. \(\lim_{t \to t_{2}} u(t) = x_{0} \pm b\) ).

\begin{enumerate}
\def\labelenumi{\arabic{enumi})}
\setcounter{enumi}{10}
\tightlist
\item
  \begin{enumerate}
  \def\labelenumii{\alph{enumii})}
  \tightlist
  \item
    On the box P: \(|t - t_{0}| < a\) , \(|x - x_{0}| < b\) let g and h be two continuous real functions such that \(g(t, x) < h(t, x)\) on P. Let u (resp. v) be an integral of \(x' = g(t, x)\) (resp. \(x' = h(t, x)\) ) such that \(u(t_{0}) = x_{0}\) (resp. \(v(t_{0}) = x_{0}\) ) defined on an interval \([t_{0}, t_{0} + c[;\) show that for \(t_{0} < t < t_{0} + c\) one has \(u(t) < v(t)\) (consider the supremum \(\tau\) of the t for which this inequality holds).
  \end{enumerate}
\end{enumerate}

\begin{enumerate}
\def\labelenumi{\alph{enumi})}
\setcounter{enumi}{1}
\item
  Let u be the maximal integral of \(x' = g(t, x)\) corresponding to the point \((t_0, x_0)\) (exerc. 10), defined on an interval \([t_0, t_0 + c[\) , with values in \(|x - x_0| < b\) . Show that in every compact interval \([t_0, t_0 + d]\) contained in \([t_0, t_0 + c[\) the minimal and maximal integrals of the equation \(x' = g(t, x) + \varepsilon\) corresponding to the point \((t_0, x_0)\) are defined once \(\varepsilon > 0\) is small enough, and converge uniformly to u when \(\varepsilon\) tends to 0 through values \textgreater{} 0.
\item
  Let g and h be two real continuous functions defined on P and such that \(g(t, x) \leqslant h(t, x)\) on P. Let \([t_0, t_0 + c[\) be an interval on which are defined an integral u of \(x' = g(t, x)\) such that \(u(t_0) = x_0\) , and the maximal integral v of \(x' = h(t, x)\) corresponding to the point \((t_0, x_0)\) . Show that \(u(t) \leqslant v(t)\) on this interval (reduce to case a) using b)).
\end{enumerate}

\begin{enumerate}
\def\labelenumi{\arabic{enumi})}
\setcounter{enumi}{11}
\tightlist
\item
  Let \(u\) be the integral of the equation \(x' = \lambda + \frac{x^2}{1 + t^2}\) equal to 0 for \(t = 0\) , and let \(J\) be the largest interval with left-hand endpoint 0 on which \(u\) is continuous.
\end{enumerate}

\begin{enumerate}
\def\labelenumi{\alph{enumi})}
\item
  Show that if \(\lambda \leqslant \frac{1}{4}\) one has \(J = [0, +\infty[\) (use exerc. 4 of IV, p.~199).
\item
  Show that if \(\lambda >\frac{1}{4}\) one has \(J = [0,a[\) where
\end{enumerate}

\[
\sinh \frac {\pi}{2 \sqrt {\lambda}} <   a <   \sinh \frac {\pi}{\sqrt {4 \lambda - 1}}
\]

(put \(x = y\sqrt{1 + t^{2}}\) and use exerc. 11).

* 13) a) Let \(\mathrm{I} = [t_0, t_0 + c]\) and let \(\omega\) be a continuous real function, \(\geqslant 0\) , defined on \(\mathrm{I} \times \mathbf{R}\) . Let \(\mathrm{S}\) be a ball with centre \(\mathbf{x}_0\) in a complete normed space \(\mathrm{E}\) , and let \(\mathbf{f}\) be a continuous map from \(\mathrm{I} \times \mathrm{S}\) into \(\mathrm{E}\) such that for any \(t \in \mathrm{I}\) , \(\mathbf{x}_1 \in \mathrm{S}\) and \(\mathbf{x}_2 \in \mathrm{S}\) one has \(\| \mathbf{f}(t, \mathbf{x}_1) - \mathbf{f}(t, \mathbf{x}_2) \| \leqslant \omega(t, \| \mathbf{x}_1 - \mathbf{x}_2 \|)\) . Let \(\mathbf{u}\) and \(\mathbf{v}\) be two integrals of \(\mathbf{x}' = \mathbf{f}(t, \mathbf{x})\) defined on \(\mathrm{I}\) , with values in \(\mathrm{S}\) , such that \(\mathbf{u}(t_0) = \mathbf{x}_1\) , \(\mathbf{v}(t_0) = \mathbf{x}_2\) ; let \(w\) be the maximal integral (exerc. 10) of \(z' = \omega(t, z)\) corresponding to the point ( \(t_0, \| \mathbf{x}_1 - \mathbf{x}_2 \|\) ), assumed to be defined on \(\mathrm{I}\) ; show that, on \(\mathrm{I}\) , one has \(\| \mathbf{u}(t) - \mathbf{v}(t) \| \leqslant w(t)\) . (Let \(w(t, \varepsilon)\) be the maximal integral of \(z' = \omega(t, z) + \varepsilon\) corresponding to the point ( \(t_0, \| \mathbf{x}_1 - \mathbf{x}_2 \|\) ), where \(\varepsilon > 0\) is sufficiently small. Show that \(\| \mathbf{u}(t) - \mathbf{v}(t) \| \leqslant w(t, \varepsilon)\) for every \(\varepsilon > 0\) , arguing by contradiction.)

\begin{enumerate}
\def\labelenumi{\alph{enumi})}
\setcounter{enumi}{1}
\tightlist
\item
  In the statement of the hypotheses of a), replace I by the interval \(I' = ]t_{0} - c, t_{0}]\) . Show that if w is the minimal integral on this interval of \(z' = \omega(t, z)\) corresponding to the point \((t_{0}, \| \mathbf{x}_{1} - \mathbf{x}_{2}\|)\) , then \(\| \mathbf{u}(t) - \mathbf{v}(t) \| \geqslant w(t)\) on \(I'\) (same method).
\end{enumerate}

* 14) a) Let \(\omega(t,z)\) be a continuous real function, \(\geqslant 0\) , defined for \(0 < t \leqslant a\) and \(z \geqslant 0\) . Suppose that \(z = 0\) is the only integral of \(z' = \omega(t,z)\) defined for \(0 < t \leqslant a\) and such that \(\lim_{t \to 0} z(t) = 0\) and \(\lim_{t \to 0} z(t)/t = 0\) . Let \(\mathrm{I} = [\mathrm{t}_0, t_0 + a[\) , let \(\mathrm{S}\) be a ball with centre \(\mathbf{x}_0\) in \(\mathrm{E}\) , and \(\mathbf{f}\) a map from \(\mathrm{I} \times \mathrm{S}\) into \(\mathrm{E}\) , such that for any \(t \in \mathrm{I}\) , \(\mathbf{x}_1 \in \mathrm{S}\) and \(\mathbf{x}_2 \in \mathrm{S}\) one has \(\|\mathbf{f}(t, \mathbf{x}_1) - \mathbf{f}(t, \mathbf{x}_2)\| \leqslant \omega(|t - t_0|, \| \mathbf{x}_1 - \mathbf{x}_2 \|)\) . Show that, on an interval with left-hand endpoint \(t_0\) contained in \(\mathrm{I}\) , the equation \(\mathbf{x}' = \mathbf{f}(t, \mathbf{x})\) can have only one solution \(\mathbf{u}\) such that \(\mathbf{u}(t_0) = \mathbf{x}_0\) . (Argue by contradiction: if \(\mathbf{v}\) is a second integral of \(\mathbf{x}' = \mathbf{f}(t, \mathbf{x})\) , bound \(\|\mathbf{u}(t) - \mathbf{v}(t)\|\) below on an interval with left-hand endpoint \(t_0\) , with the help of exerc. 13 b), and thus obtain a contradiction.)

Apply to the case where \(\omega(t, z) = k(z/t)\) with \(0 \leqslant k < 1\) (cf.~IV, p.~200, exerc. 8).

\begin{enumerate}
\def\labelenumi{\alph{enumi})}
\setcounter{enumi}{1}
\item
  The result of a) applies for \(\omega(t,z)=z/t\) ; but show by an example in this case that if u, v are two integrals to within \(\varepsilon\) of \(\mathbf{x}^{\prime}=\mathbf{f}(t,\mathbf{x})\) , equal to \(x_{0}\) at the point \(t_{0}\) , it is not possible to bound \(\|\mathbf{u}(t)-\mathbf{v}(t)\|\) above by a number depending only on t (and not on the function f). (Take for f a continuous map of \(R_{+}\times R\) into R, equal to x/t for \(t\geqslant\alpha\) and for \(0<t<\alpha\) and \(|x|\leqslant t^{2}/(\alpha-t)\) , and independent of x for the other points \((t,x)\) .)
\item
  Let \(\theta\) be a continuous real function, \(\geqslant 0\) on the interval \([0, a]\) . Show that if the integral \(\int_0^a \frac{\theta(t)}{t} dt\) converges then the result of \(a)\) applies for \(\omega(t, z) = \frac{1 + \theta(t)}{t} z\) ; on the other hand, if \(\int_0^a \frac{\theta(t)}{t} dt\) is infinite, give an example of a continuous function \(\mathbf{f}\) such that \(\| \mathbf{f}(t, \mathbf{x}_1) - \mathbf{f}(t, \mathbf{x}_2) \| \leqslant \frac{1 + \theta(|t - t_0|)}{|t - t_0|} \| \mathbf{x}_1 - \mathbf{x}_2 \|\) , but such that the equation \(\mathbf{x}' = \mathbf{f}(t, \mathbf{x})\) has infinitely many integrals equal to \(\mathbf{x}_0\) at the point \(t_0\) (method similar to that of \(b\) )).
\end{enumerate}

\begin{enumerate}
\def\labelenumi{\arabic{enumi})}
\setcounter{enumi}{14}
\item
  Let \(f\) be a real function, defined and continuous for \(|t| \leqslant a\) , \(|x| \leqslant b\) , such that \(f(t,x) < 0\) for \(tx > 0\) and \(f(t,x) > 0\) for \(tx < 0\) ; show that \(x = 0\) is the only integral of the equation \(x' = f(t,x)\) which takes the value 0 at the point \(t = 0\) (argue by contradiction).
\item
  Let E be a vector space of finite dimension, and f a bounded function on I×H satisfying the conditions of IV, p.~165, lemma 1, such that the equation \(\mathbf{x}^{\prime}=\mathbf{f}(t,\mathbf{x})\) has a unique solution u defined on I, with values in H, equal to \(x_{0}\) at the point \(t_{0}\) . Suppose further that for sufficiently large n there is an approximate integral \(u_{n}\) of \(\mathbf{x}^{\prime}=\mathbf{f}(t,\mathbf{x})\) to within 1/n, defined on I, with values in H, equal to \(x_{0}\) at the point \(t_{0}\) . Show that the sequence \((\mathbf{u}_{n})\) converges uniformly to u on every compact interval contained in I (use the fact that the sequence \((\mathbf{u}_{n})\) is equicontinuous on I).
\item
  To study the equation \(x' = \sin tx\) (IV, p.~174, Example 4) one puts u = xy and considers the solutions of the corresponding equation \(u' = \frac{u}{t} + t \sin u = \mathrm{F}(t, u)\) which are \textgreater0 for t \textgreater0 (one has \(u(0) = 0\) for every solution of this kind). Denote by \(\Gamma_k\) the curve defined by the relations \((2k - 1)\pi < u < 2k\pi\) , \(t \sin u + \frac{u}{t} = 0\) for each integer \(k \geqslant 1\) ; by \(D_k\) the open set defined by \((2k - 1)\pi < u < 2k\pi\) , \(t \sin u + \frac{u}{t} < 0\) ; and by E the complement of the union of the \(\overline{D}_k\) in the set of \((t, u)\) such that t \textgreater0 and u \textgreater0.
\end{enumerate}

\begin{enumerate}
\def\labelenumi{\alph{enumi})}
\item
  Show that every integral curve which cuts a line \(u = 2k\pi\) also cuts the line \(u = (2k + 1)\pi\) .
\item
  If an integral curve C cuts a curve \(\Gamma_{k}\) at a point \((t_{0}, u_{0})\) , then the function u is increasing for 0 \textless{} t \textless{} t \textless{} \(t_{0}\), decreasing for \(t > t_{0}\) , and, when t tends to \(+\infty\) , \(u(t)\) tends to \((2k - 1)\pi\) , the curve C remaining in \(D_{k}\) for \(t > t_{0}\) .
\item
  Show that there is no integral curve C contained in E and such that \(u(t)\) tends to \(+\infty\) when t tends to \(+\infty\) . (Form the differential equation dx/du = G(u, x) between x and u along C. If C lies entirely in E, then \(x^{2} > u\) for \(u = (2k - \frac{1}{2})\pi\) ; on the other hand, estimate \(x^{2}(u)\) by using the preceding differential equation, and obtain a contradiction.)
\end{enumerate}

\(d\)) Show that for every integer \(k\) there exists one and only one integral curve C contained in E and such that \(u(t)\) tends to \(2k\pi\) when \(t\) tends to \(+\infty\) . (On putting \(v = 2k\pi - u\) one obtains the equation \(v' = \frac{v - 2k\pi}{t} + t \sin v\) ; compare this equation to \(v' = \frac{v - 2k\pi}{t} + tv\) with \(t\) tending to \(+\infty\) and \(v\) to 0.)

\begin{enumerate}
\def\labelenumi{\arabic{enumi})}
\setcounter{enumi}{17}
\tightlist
\item
  Let E be the space of sequences \(\mathbf{x}=(x_{n})_{n\in\mathbb{N}}\) of real numbers such that \(\lim_{n\to\infty}x_{n}=0\) , endowed with the norm \(\|x\|=\sup_{n}|x_{n}|\) , which is in fact a complete normed space. For each \(\mathbf{x}=(x_{n})\in\mathrm{E}\) let \(\mathbf{y}=\mathbf{f}(\mathbf{x})\) be the element \((y_{n})\) of E defined by \(y_{n}=|x_{n}|^{\frac{1}{2}}+\frac{1}{n+1}\) ; the function f is continuous on E. Show that there is no solution of the differential equation \(\mathbf{x}^{\prime}=\mathbf{f}(\mathbf{x})\) defined on a neighbourhood of 0 in R, with values in E, and equal to 0 at the point t=0 (cf.~IV, p.~202, exerc. 11).
\end{enumerate}

\setcounter{exercisegroup}{1}
\refstepcounter{exercisegroup}
